\section{An existence theorem for coherent algebraic sheaves}
\label{section:III.5}


\subsection{Statement of the theorem}
\label{subsection:III.5.1}

\begin{env}[5.1.1]
\label{III.5.1.1}
Let $A$ be an adic Noetherian ring and $\mathfrak{J}$ an ideal of
definition of $A$, so that $A$ is separated and complete for the
$\mathfrak{J}$-adic topology.  If $Y=\Spec(A)$, the affine formal
prescheme $\Spf(A)$ is identified with the completion
$\widehat{Y}$ of $Y$ along the closed subset
$Y'=\mathrm{V}(\mathfrak{J})$ \sref[I]{I.10.10.1}.  Let $X$ be an
ordinary $Y$-prescheme of finite type, with structural morphism
$f:X\to Y$.  We denote by $\widehat{X}$ the completion of $X$ along
the closed subset $X'=f^{-1}(Y')$, or equivalently the formal
$\widehat{Y}$-prescheme $X\times_Y\widehat{Y}$, and by
$\widehat{f}:\widehat{X}\to\widehat{Y}$ the extension of $f$ to the
completions.  Finally, for every coherent $\sh{O}_X$-module
$\sh{F}$, we denote by $\widehat{\sh{F}}$ its completion
$\sh{F}_{/X'}$, which is a coherent
$\sh{O}_{\widehat{X}}$-module.
\end{env}

\begin{proposition}[5.1.2]
\label{III.5.1.2}
With the hypotheses and notation of \sref{III.5.1.1}, let $\sh{F}$
be a coherent $\sh{O}_X$-module whose support is proper over $Y$
\sref[II]{II.5.4.10}.  The canonical homomorphisms
\sref{III.4.1.4}
\[
  \rho_i:
  \HH^i(X,\sh{F})
  \longrightarrow
  \HH^i(\widehat{X},\widehat{\sh{F}})
\]
are isomorphisms.
\end{proposition}

\begin{proof}
The $A$-module $\HH^i(X,\sh{F})$ is of finite type
\sref{III.3.2.4}, and is therefore identical with its separated
completion for the $\mathfrak{J}$-preadic topology
\sref[0]{0.7.3.6}.  The proposition is consequently a special case
of \sref{III.4.1.10}.
\end{proof}

Recall that the canonical isomorphisms $\rho_i$ commute with the
connecting homomorphisms associated with every exact sequence
\[
  0\longrightarrow\sh{F}'
  \longrightarrow\sh{F}
  \longrightarrow\sh{F}''
  \longrightarrow0
\]
of coherent $\sh{O}_X$-modules
\sref[0]{0.12.1.6} \sref[I]{I.10.8.9}.

\begin{corollary}[5.1.3]
\label{III.5.1.3}
Let $\sh{F}$ and $\sh{G}$ be coherent $\sh{O}_X$-modules such that
the intersection of their supports is proper over $Y$.  The
canonical homomorphism
\[
  \Hom_{\sh{O}_X}(\sh{F},\sh{G})
  \longrightarrow
  \Hom_{\sh{O}_{\widehat{X}}}
  (\widehat{\sh{F}},\widehat{\sh{G}})
  \tag{5.1.3.1}
  \label{III.5.1.3.1.eq}
\]
which sends a homomorphism $u:\sh{F}\to\sh{G}$ to its completion
$\widehat{u}:\widehat{\sh{F}}\to\widehat{\sh{G}}$, is an
isomorphism.  Moreover, when $f$ is closed,
$\widehat{u}$ is injective (respectively, surjective) if and only
if $u$ is injective (respectively, surjective).
\end{corollary}

\begin{proof}
The first assertion is a special case of \sref{III.4.5.3}; it also
follows from the fact that the source of \eqref{III.5.1.3.1.eq} is
an $A$-module of finite type and is therefore identical with its
separated completion.  For the second assertion, note that, by
\sref[I]{I.10.8.14}, $\widehat{u}$ is injective (respectively,
surjective) if and only if there is a neighbourhood of $X'$ in
$X$ on which $u$ is injective (respectively, surjective).  The
conclusion therefore follows from the lemma below.
\end{proof}

\oldpage[III]{150}

\begin{lemma}[5.1.3.1]
\label{III.5.1.3.1}
Under the hypotheses of \sref{III.5.1.1}, suppose in addition that
$f$ is closed.  Every neighbourhood of $X'$ in $X$ is then equal
to $X$.
\end{lemma}

\begin{proof}
We first reduce to the case $f(X)=Y$.  Indeed, by hypothesis,
$f(X)$ is a closed subset $Z$ of $Y$.  We may replace $f$ by
$f_{\mathrm{red}}$ \sref[I]{I.6.3.4}, and hence suppose that $X$
and $Y$ are reduced; we may then replace $Y$ by the reduced closed
subprescheme having $Z$ as its underlying space
\sref[I]{I.5.2.2}, since every ideal of $A$ is closed and every
quotient ring of $A$ is therefore adic and Noetherian.

We then have $f(X')=Y'$.  If $V$ is an open neighbourhood of $X'$
in $X$, the set $f(X-V)$ is closed in $Y$ and does not meet $Y'$.
This is impossible unless $X-V$ is empty, because $\mathfrak{J}$
is contained in the radical of $A$
\sref[I]{I.1.1.15} \sref[0]{0.7.1.10}.  Thus $V=X$.
\end{proof}

If we restrict attention to coherent $\sh{O}_X$-modules whose
support is proper over $Y$, \sref{III.5.1.3} says, in categorical
language, that the functor
\[
  \sh{F}\longmapsto\widehat{\sh{F}}
\]
is fully faithful from the category of such $\sh{O}_X$-modules to
the category of coherent $\sh{O}_{\widehat{X}}$-modules.  It thus
gives an equivalence of the former category with a full
subcategory of the latter \sref[0]{0.8.1.6}.  The existence theorem
will show that, when $X$ is proper over $Y$, this subcategory is in
fact the category of all coherent
$\sh{O}_{\widehat{X}}$-modules.  More precisely:

\begin{theorem}[5.1.4]
\label{III.5.1.4}
Let $A$ be an adic Noetherian ring, let $Y=\Spec(A)$, let
$\mathfrak{J}$ be an ideal of definition of $A$, and put
$Y'=\mathrm{V}(\mathfrak{J})$.  Let $f:X\to Y$ be a separated
morphism of finite type and put $X'=f^{-1}(Y')$.  Let
\[
  \widehat{Y}=Y_{/Y'}=\Spf(A),\qquad
  \widehat{X}=X_{/X'}
\]
be the completions of $Y$ and $X$ along $Y'$ and $X'$, and let
$\widehat{f}:\widehat{X}\to\widehat{Y}$ be the extension of $f$ to
the completions.  Then the functor
\[
  \sh{F}\longmapsto\sh{F}_{/X'}=\widehat{\sh{F}}
\]
is an equivalence from the category of coherent
$\sh{O}_X$-modules whose support is proper over $\Spec(A)$ to the
category of coherent $\sh{O}_{\widehat{X}}$-modules whose support
is proper over $\Spf(A)$.
\end{theorem}

In other words, in view of \sref{III.5.1.3}:

\begin{corollary}[5.1.5]
\label{III.5.1.5}
An $\sh{O}_{\widehat{X}}$-module is isomorphic to the completion of
a coherent $\sh{O}_X$-module whose support is proper over
$\Spec(A)$ if and only if it is coherent and its support is proper
over $\Spf(A)$.
\end{corollary}

The most important case is the following:

\begin{corollary}[5.1.6]
\label{III.5.1.6}
Suppose that $X$ is proper over $Y=\Spec(A)$.  Then the functor
$\sh{F}\mapsto\widehat{\sh{F}}$ is an equivalence between the
category of coherent $\sh{O}_X$-modules and the category of
coherent $\sh{O}_{\widehat{X}}$-modules.
\end{corollary}

\begin{env}[5.1.7]
\label{III.5.1.7}
\emph{Scholium.}
Taking account of the characterization of coherent sheaves on
formal preschemes \sref[I]{I.10.11.3}, we see that, under the
conditions of \sref{III.5.1.1}, giving a coherent
$\sh{O}_X$-module whose support is proper over $\Spec(A)$ is
equivalent---if
\[
  Y_n=\Spec(A/\mathfrak{J}^{n+1}),\qquad
  X_n=X\times_Y Y_n
\]
---to giving a projective system $(\sh{F}_n)$ of coherent
$\sh{O}_{X_n}$-modules such that, for $m\leq n$,
\[
  \sh{F}_m
  =
  \sh{F}_n\otimes_{\sh{O}_{Y_n}}\sh{O}_{Y_m},
  \qquad\text{or equivalently}\qquad
  \sh{F}_m
  =
  \sh{F}_n/\mathfrak{J}^{m+1}\sh{F}_n,
\]
and such that the support of $\sh{F}_0$ is a subset of $X_0$
proper over $Y_0$.  By \sref[I]{I.10.11.4}, homomorphisms of
coherent $\sh{O}_X$-modules are interpreted similarly as
homomorphisms of projective systems of coherent
$\sh{O}_{X_n}$-modules.

\oldpage[III]{151}

In all known applications, $A$ is in fact a local adic Noetherian
ring.  Thus the $Y_n$ are spectra of local Artinian rings, and the
results of this subsection and of the preceding ones reduce, to a
large extent, algebraic geometry over a local adic Noetherian ring
to algebraic geometry over local Artinian rings.
\end{env}

\begin{corollary}[5.1.8]
\label{III.5.1.8}
Under the hypotheses of \sref{III.5.1.4}, the map
\[
  Z\longmapsto\widehat{Z}=Z_{/(Z\cap X')}
\]
is a bijection from the set of closed subpreschemes $Z$ of $X$
which are proper over $Y$ onto the set of closed formal
subpreschemes of $\widehat{X}$ which are proper over
$\widehat{Y}$.
\end{corollary}

\begin{proof}
A closed formal subprescheme of $\widehat{X}$ has the form
\[
  \bigl(T,
    (\sh{O}_{\widehat{X}}/\sh{A})|_T\bigr),
\]
where $\sh{A}$ is a coherent ideal of
$\sh{O}_{\widehat{X}}$ \sref[I]{I.10.14.2}.  If $T$ is proper over
$\widehat{Y}$, \sref{III.5.1.4} shows that
$\sh{O}_{\widehat{X}}/\sh{A}$ is isomorphic to an
$\sh{O}_{\widehat{X}}$-module of the form
$\widehat{\sh{F}}$, where $\sh{F}$ is a coherent
$\sh{O}_X$-module whose support is proper over $Y$.  Moreover,
\sref{III.5.1.3} shows that the canonical homomorphism
\[
  \sh{O}_{\widehat{X}}
  \longrightarrow
  \sh{O}_{\widehat{X}}/\sh{A}
\]
has the form $\widehat{u}$, where
$u:\sh{O}_X\to\sh{F}$ is a surjective homomorphism of
$\sh{O}_X$-modules.  Consequently
$\sh{F}=\sh{O}_X/\sh{N}$ for a coherent ideal $\sh{N}$ of
$\sh{O}_X$, and
$\sh{A}=\widehat{\sh{N}}$ \sref[I]{I.10.8.8}.  The conclusion now
follows from \sref[I]{I.10.14.7}.
\end{proof}


\subsection{Proof of the existence theorem: the projective and quasi-projective case}
\label{subsection:III.5.2}

\begin{env}[5.2.1]
\label{III.5.2.1}
Under the conditions of \sref{III.5.1.4}, we shall say
provisionally that a coherent $\sh{O}_{\widehat{X}}$-module is
\emph{algebraizable} if it is isomorphic to the completion
$\widehat{\sh{F}}$ of a coherent $\sh{O}_X$-module $\sh{F}$ whose
support is proper over $Y$.
\end{env}

\begin{lemma}[5.2.2]
\label{III.5.2.2}
Let $\sh{F}'$ and $\sh{G}'$ be algebraizable
$\sh{O}_{\widehat{X}}$-modules.  For every homomorphism
$u:\sh{F}'\to\sh{G}'$, the modules
$\Ker(u)$, $\Im(u)$, and $\Coker(u)$ are algebraizable.
\end{lemma}

\begin{proof}
Write
\[
  \sh{F}'=\widehat{\sh{F}},\qquad
  \sh{G}'=\widehat{\sh{G}},
\]
where $\sh{F}$ and $\sh{G}$ are coherent $\sh{O}_X$-modules whose
supports are proper over $Y$.  By \sref{III.5.1.3},
$u=\widehat{v}$ for a homomorphism
$v:\sh{F}\to\sh{G}$.  Since completion is exact,
$\Ker(\widehat{v})\simeq\widehat{\Ker(v)}$; and because the support
of $\Ker(v)$ is contained in that of $\sh{F}$, it follows that
$\Ker(u)$ is algebraizable.  The arguments for $\Im(u)$ and
$\Coker(u)$ are the same.
\end{proof}

\begin{proposition}[5.2.3]
\label{III.5.2.3}
Let $A$ be an adic Noetherian ring, $\mathfrak{J}$ an ideal of
definition of $A$, $\mathfrak{Y}=\Spf(A)$, and
$f:\mathfrak{X}\to\mathfrak{Y}$ a proper morphism of formal
preschemes.  Put
\[
  Y_k=\Spec(A/\mathfrak{J}^{k+1}),\qquad
  X_k=\mathfrak{X}\times_{\mathfrak{Y}}Y_k,
\]
and, for every $\sh{O}_{\mathfrak{X}}$-module $\sh{F}$, put
\[
  \sh{F}_k
  =
  \sh{F}\otimes_{\sh{O}_{\mathfrak{X}}}\sh{O}_{X_k}
  =
  \sh{F}/\mathfrak{J}^{k+1}\sh{F}.
\]
Let $\sh{L}$ be an invertible $\sh{O}_{\mathfrak{X}}$-module, and
suppose that
$\sh{L}_0=\sh{L}/\mathfrak{J}\sh{L}$ is an ample
$\sh{O}_{X_0}$-module.  For every
$\sh{O}_{\mathfrak{X}}$-module $\sh{F}$ and every integer $n$, set
\[
  \sh{F}(n)=\sh{F}\otimes\sh{L}^{\otimes n}.
\]
Then, for every coherent $\sh{O}_{\mathfrak{X}}$-module $\sh{F}$,
there is an integer $n_0$ such that, for every $n\geq n_0$:
\begin{enumerate}
  \item[\rm{(i)}] the canonical homomorphism
    \[
      \HH^0(\mathfrak{X},\sh{F}(n))
      \longrightarrow
      \HH^0(\mathfrak{X},\sh{F}_k(n))
    \]
    is surjective for every $k\geq0$;
  \item[\rm{(ii)}]
    $\HH^q(\mathfrak{X},\sh{F}(n))=0$ for every $q>0$.
\end{enumerate}
\end{proposition}

\begin{proof}
The underlying spaces of $\mathfrak{X}$ and $X_0$ are the same.
The sheaves
\[
  \sh{M}_k
  =
  \mathfrak{J}^k\sh{F}/\mathfrak{J}^{k+1}\sh{F},
\]
being annihilated by $\mathfrak{J}$, may be regarded as coherent
$\sh{O}_{X_0}$-modules \sref[0]{0.5.3.10}.  Moreover, if
\[
  \sh{M}_k(n)
  =
  \sh{M}_k\otimes_{\sh{O}_{X_0}}\sh{L}_0^{\otimes n},
\]
then
\[
  \sh{M}_k(n)
  =
  \mathfrak{J}^k\sh{F}(n)/
  \mathfrak{J}^{k+1}\sh{F}(n).
\]
Since $\sh{L}_0$ is ample for $f_0:X_0\to Y_0$ and $f_0$ is
proper,
\oldpage[III]{152}
$f_0$ is projective \sref[II]{II.5.5.4}.  Let
\[
  S=\bigoplus_{k\geq0}
  \mathfrak{J}^k/\mathfrak{J}^{k+1}
\]
be the graded $A$-algebra associated with the
$\mathfrak{J}$-adic filtration of $A$.  It is of finite type
because $A$ is Noetherian.  Put
\[
  \sh{S}'=f_0^*(\widetilde{S}),\qquad
  \sh{M}=\bigoplus_{k\geq0}\sh{M}_k.
\]
Then $\sh{S}'$ is a quasi-coherent $\sh{O}_{X_0}$-algebra of
finite type, and $\sh{M}$ is a quasi-coherent graded
$\sh{S}'$-module of finite type, since $\sh{F}_0$ is coherent and
generates $\sh{M}$ as an $\sh{S}'$-module.  Theorem
\sref{III.2.4.1},~\emph{(ii)}, therefore gives an integer $n_0$
such that, for $n\geq n_0$ and every $k$,
\[
  \HH^q(X_0,\sh{M}_k(n))=0
  \qquad(q>0).
  \tag{5.2.3.1}
  \label{III.5.2.3.1}
\]
Regarded now as an $\sh{O}_{\mathfrak{X}}$-module,
$\sh{M}_k(n)$ consequently also satisfies
\[
  \HH^q(\mathfrak{X},\sh{M}_k(n))=0
  \qquad(q>0,\ n\geq n_0).
\]

Apply the exact cohomology sequence to
\[
  0\longrightarrow
  \frac{\mathfrak{J}^k\sh{F}(n)}
       {\mathfrak{J}^{k+1}\sh{F}(n)}
  \longrightarrow
  \frac{\mathfrak{J}^h\sh{F}(n)}
       {\mathfrak{J}^{k+1}\sh{F}(n)}
  \longrightarrow
  \frac{\mathfrak{J}^h\sh{F}(n)}
       {\mathfrak{J}^{k}\sh{F}(n)}
  \longrightarrow0.
\]
For $0\leq h<k$, $n\geq n_0$, and $q>0$, induction on $k-h$
gives
\[
  \HH^q\left(
    \mathfrak{X},
    \frac{\mathfrak{J}^h\sh{F}(n)}
         {\mathfrak{J}^{k}\sh{F}(n)}
  \right)=0,
  \tag{5.2.3.2}
  \label{III.5.2.3.2}
\]
and in particular, for $h=0$,
\[
  \HH^q(\mathfrak{X},\sh{F}_k(n))=0
  \qquad(n\geq n_0,\ k\geq0,\ q>0).
  \tag{5.2.3.3}
  \label{III.5.2.3.3}
\]
Another part of the exact cohomology sequence, with $h=0$, gives
\[
  \HH^0(\mathfrak{X},\sh{F}_{k+1}(n))
  \longrightarrow
  \HH^0(\mathfrak{X},\sh{F}_k(n))
  \longrightarrow
  \HH^1\left(
    \mathfrak{X},
    \frac{\mathfrak{J}^k\sh{F}(n)}
         {\mathfrak{J}^{k+1}\sh{F}(n)}
  \right)
  =0.
  \tag{5.2.3.4}
  \label{III.5.2.3.4}
\]
It follows that, for $h\leq k$, the canonical map
\[
  \HH^0(\mathfrak{X},\sh{F}_k(n))
  \longrightarrow
  \HH^0(\mathfrak{X},\sh{F}_h(n))
  \tag{5.2.3.5}
  \label{III.5.2.3.5}
\]
is surjective.  Thus, for every $q$, the projective system
\[
  \bigl(\HH^q(\mathfrak{X},\sh{F}_k(n))\bigr)_{k\geq0}
\]
satisfies condition (ML) when $n\geq n_0$.

Furthermore, every formal affine open $U$ of $\mathfrak{X}$ is an
affine open in each $X_k$ \sref[I]{I.10.5.2}.  Hence
$\HH^q(U,\sh{F}_k(n))=0$ for $q>0$ \sref[I]{I.1.3.1}, and
\[
  \HH^0(U,\sh{F}_k(n))
  \longrightarrow
  \HH^0(U,\sh{F}_h(n))
\]
is surjective for $h\leq k$ \sref[I]{I.1.3.9}.  The hypotheses of
\sref[0]{0.13.3.1} are therefore satisfied.  For $n\geq n_0$ we
obtain:
\begin{enumerate}
  \item[$1^\circ$] for every $q>0$, the map
    \[
      \HH^q(\mathfrak{X},\sh{F}(n))
      \longrightarrow
      \varprojlim_k
      \HH^q(\mathfrak{X},\sh{F}_k(n))
    \]
    is bijective; hence, by \eqref{III.5.2.3.3},
    $\HH^q(\mathfrak{X},\sh{F}(n))=0$;
  \item[$2^\circ$] the homomorphism
    \[
      \HH^0(\mathfrak{X},\sh{F}(n))
      \longrightarrow
      \varprojlim_k
      \HH^0(\mathfrak{X},\sh{F}_k(n))
    \]
    is bijective.  Since the homomorphisms
    \eqref{III.5.2.3.5} are surjective, so is each homomorphism
    \[
      \varprojlim_k
      \HH^0(\mathfrak{X},\sh{F}_k(n))
      \longrightarrow
      \HH^0(\mathfrak{X},\sh{F}_h(n)).
    \]
\end{enumerate}
This proves both assertions.
\end{proof}

\begin{corollary}[5.2.4]
\label{III.5.2.4}
Under the hypotheses of \sref{III.5.2.3}, for every coherent
$\sh{O}_{\mathfrak{X}}$-module $\sh{F}$ there is an integer $N$
such that, for $n\geq N$, $\sh{F}(n)$ is generated by its sections
over
\oldpage[III]{153}
$\mathfrak{X}$.  Equivalently, $\sh{F}$ is isomorphic to a quotient
of an $\sh{O}_{\mathfrak{X}}$-module of the form
$(\sh{O}_{\mathfrak{X}}(-n))^k$.
\end{corollary}

\begin{proof}
Since $X_0$ is Noetherian, the hypothesis on $\sh{L}_0$ and
\sref[II]{II.4.5.5} give an integer $n_0$ such that
$\sh{F}_0(n)$ is generated by its sections over $\mathfrak{X}$ for
$n\geq n_0$.  We may take $n_0$ large enough that
\[
  \Gamma(\mathfrak{X},\sh{F}(n))
  \longrightarrow
  \Gamma(\mathfrak{X},\sh{F}_0(n))
\]
is also surjective for $n\geq n_0$ \sref{III.5.2.3}.  There are
therefore finitely many sections
$s_i\in\Gamma(\mathfrak{X},\sh{F}(n))$ whose images generate
$\sh{F}_0(n)$ \sref[0]{0.5.2.3}.  Since $\mathfrak{J}$ is
contained in the maximal ideal of the local ring at every point
of $\mathfrak{X}$, Nakayama's lemma applied to those local rings
shows that the $s_i$ generate $\sh{F}(n)$
\sref[0]{0.5.1.1}.
\end{proof}

\begin{env}[5.2.5]
\label{III.5.2.5}
\emph{Proof of the existence theorem: the projective case.}
With the notation of \sref{III.5.1.4}, suppose that $f$ is
projective, so that there is an ample $\sh{O}_X$-module $\sh{L}$
\sref[I]{I.5.5.4}.  By definition,
\[
  X_n=\widehat{X}\times_{\widehat{Y}}Y_n
\]
is the closed subprescheme
$X\times_Y Y_n=f^{-1}(Y_n)$ of $X$.  If
\[
  \sh{L}'=\sh{L}\otimes_{\sh{O}_X}\sh{O}_{\widehat{X}}
\]
is the completion of $\sh{L}$, then
$\sh{L}'_0=\sh{L}/\mathfrak{J}\sh{L}$, regarded as an
$\sh{O}_{X_0}$-module; this module is ample
\sref[II]{II.4.6.13}[(i~\emph{bis})].

We may therefore apply \sref{III.5.2.4} to $\sh{L}'$ and to any
coherent $\sh{O}_{\widehat{X}}$-module $\sh{F}$.  It follows that
$\sh{F}$ is isomorphic to a quotient of
\[
  \sh{G}=
  \bigl(\sh{L}'^{\otimes(-n)}\bigr)^k
\]
for suitable integers $n>0$ and $k>0$.  The module $\sh{G}$ is the
completion of $(\sh{L}^{\otimes(-n)})^k$
\sref[I]{I.10.8.10}, and is therefore algebraizable.  Let
$u:\sh{G}\to\sh{F}$ be the canonical surjection, and put
$\sh{H}=\Ker(u)$; this is a coherent
$\sh{O}_{\widehat{X}}$-module \sref[0]{0.5.3.4}.  Applying the same
argument to $\sh{H}$, we find an algebraizable
$\sh{O}_{\widehat{X}}$-module $\sh{K}$ and a homomorphism
$v:\sh{K}\to\sh{G}$ such that $\sh{H}=\Im(v)$.  Thus
$\sh{F}=\Coker(v)$, and $\sh{F}$ is algebraizable by
\sref{III.5.2.2}.
\end{env}

\begin{env}[5.2.6]
\label{III.5.2.6}
\emph{Proof of the existence theorem: the quasi-projective case.}
Retain the notation of \sref{III.5.1.4}, and now suppose that $f$
is quasi-projective.  There is a projective morphism
$g:Z\to Y$ such that $X$ is identified with the $Y$-prescheme
induced on an open subset of $Z$ \sref[II]{II.5.3.2}.  If
$Z'=g^{-1}(Y')$, then $X'=X\cap Z'$.  Consequently the completion
$\widehat{X}=X_{/X'}$ is identified with the formal prescheme
induced by $\widehat{Z}=Z_{/Z'}$ on the open subset $X\cap Z'$ of
$\widehat{Z}$ \sref[I]{I.10.8.5}.

Let $\sh{F}'$ be a coherent $\sh{O}_{\widehat{X}}$-module whose
support $T'$ is proper over $\widehat{Y}$.  By definition, there
is a closed subprescheme of $X'$ with underlying space
$T'\subset X'$ such that the restriction $T'\to Y'$ of $f$ is
proper.  Hence $T'$ is proper over $Y$, and is therefore closed in
$Z'$ \sref[II]{II.5.4.10}.  It follows that $\sh{F}'$ is induced
on $\widehat{X}$ by the $\sh{O}_{\widehat{Z}}$-module $\sh{G}'$
obtained by gluing $\sh{F}'$, on the open $\widehat{X}$, with the
zero sheaf on the open $\widehat{Z}-T'$; the two sheaves agree on
their intersection $\widehat{X}-T'$.

The module $\sh{G}'$ is coherent.  By \sref{III.5.2.5}, there is a
coherent $\sh{O}_Z$-module $\sh{G}$ such that
$\sh{G}'=\widehat{\sh{G}}$.  Let $T$ be the support of $\sh{G}$;
then $T'=T\cap Z'$ \sref[I]{I.10.8.12}.  If $h$ is the restriction
of $g$ to the reduced closed subprescheme of $Z$ having $T$ as
its underlying space, then
\[
  T'=h^{-1}(Y')=T\cap g^{-1}(Y').
\]
Thus $X\cap T$ is an open subset of $T$ containing $T'$.

\oldpage[III]{154}
Since $h$ is proper \sref[II]{II.5.4.2}, and hence closed,
\sref{III.5.1.3.1} gives $X\cap T=T$.  Thus $T\subset X$; and
since $T$ is closed in $Z$, it is proper over $Y$.  If $\sh{F}$ is
the sheaf induced on $X$ by $\sh{G}$, its completion
$\widehat{\sh{F}}$ is induced on $\widehat{X}$ by
$\widehat{\sh{G}}$ \sref[I]{I.10.8.4}, and is therefore equal to
$\sh{F}'$.  This completes the proof.
\end{env}


\subsection{Proof of the existence theorem: the general case}
\label{subsection:III.5.3}

\begin{lemma}[5.3.1]
\label{III.5.3.1}
Under the conditions of \sref{III.5.1.4}, let
\[
  0\longrightarrow\sh{H}\longrightarrow\sh{F}
  \longrightarrow\sh{G}\longrightarrow0
\]
be an exact sequence of coherent
$\sh{O}_{\widehat{X}}$-modules.  If $\sh{G}$ and $\sh{H}$ are
algebraizable, then $\sh{F}$ is algebraizable.
\end{lemma}

\begin{proof}
Suppose that
\[
  \sh{G}=\widehat{\sh{B}},
  \qquad
  \sh{H}=\widehat{\sh{C}},
\]
where $\sh{B}$ and $\sh{C}$ are coherent $\sh{O}_X$-modules whose
supports are proper over $Y$.  The module $\sh{F}$ canonically
defines an element of the $A$-module
\[
  \Ext^1_{\sh{O}_{\widehat{X}}}
  (\widehat{X};\widehat{\sh{B}},\widehat{\sh{C}})
\]
\sref[0]{0.12.3.2}.  The hypotheses imply that this $A$-module is
canonically isomorphic to
\[
  \Ext^1_{\sh{O}_X}(X;\sh{B},\sh{C})
\]
\sref{III.4.5.2}.  There is therefore an exact sequence
\[
  0\longrightarrow\sh{C}\longrightarrow\sh{A}
  \longrightarrow\sh{B}\longrightarrow0
\]
of coherent $\sh{O}_X$-modules whose canonical class maps to the
class corresponding to $\sh{F}$.  By definition, taking account of
\sref[I]{I.10.8.8}, (ii), this says that $\sh{F}$ is isomorphic to
$\widehat{\sh{A}}$.  The lemma follows, since
$\Supp(\sh{A})$ is contained in
$\Supp(\sh{B})\cup\Supp(\sh{C})$ and is therefore proper over $Y$.
\end{proof}

\begin{corollary}[5.3.2]
\label{III.5.3.2}
Under the conditions of \sref{III.5.1.1}, let
$u:\sh{F}\to\sh{G}$ be a homomorphism of coherent
$\sh{O}_{\widehat{X}}$-modules.  If $\sh{G}$, $\Ker(u)$, and
$\Coker(u)$ are algebraizable, then $\sh{F}$ is algebraizable.
\end{corollary}

\begin{proof}
Lemma \sref{III.5.2.2}, applied to
$\sh{G}\to\Coker(u)$, shows that $\Im(u)$ is algebraizable.  It then
suffices to apply Lemma \sref{III.5.3.1} to the exact sequence
\[
  0\longrightarrow\Ker(u)\longrightarrow\sh{F}
  \longrightarrow\Im(u)\longrightarrow0.
\]
\end{proof}

\begin{lemma}[5.3.3]
\label{III.5.3.3}
Under the conditions of \sref{III.5.1.1}, let $h:Z\to Y$ be a
morphism of finite type, let $\widehat{Z}$ be the completion of $Z$
along $Z'=h^{-1}(Y')$, let $g:Z\to X$ be a proper $Y$-morphism,
and let $\widehat{g}:\widehat{Z}\to\widehat{X}$ be its extension to
the completions.  For every algebraizable
$\sh{O}_{\widehat{Z}}$-module $\sh{F}'$, the module
$\widehat{g}_*(\sh{F}')$ is an algebraizable
$\sh{O}_{\widehat{X}}$-module.
\end{lemma}

\begin{proof}
If $\sh{F}$ is a coherent $\sh{O}_Z$-module such that
$\sh{F}'=\widehat{\sh{F}}$, the first comparison theorem
\sref{III.4.1.5} identifies
$\widehat{g}_*(\widehat{\sh{F}})$ with the completion of
$g_*(\sh{F})$.
\end{proof}

\begin{lemma}[5.3.4]
\label{III.5.3.4}
Let $X$ be an ordinary Noetherian scheme, let $X'$ be a closed
subset of $X$, let $f:Z\to X$ be a proper morphism, and put
\[
  Z'=f^{-1}(X'),\qquad
  \widehat{X}=X_{/X'},\qquad
  \widehat{Z}=Z_{/Z'}.
\]
Let $\widehat{f}:\widehat{Z}\to\widehat{X}$ be the extension of $f$
to the completions.  Let $\sh{M}$ be a coherent ideal of
$\sh{O}_X$ such that, if
\[
  U=X-\Supp(\sh{O}_X/\sh{M}),
\]
the restriction $f^{-1}(U)\to U$ of $f$ is an isomorphism.  Then,
for every coherent $\sh{O}_{\widehat{X}}$-module $\sh{F}$, there is
an integer $n>0$ such that the kernel and cokernel of the canonical
homomorphism
\[
  \rho_{\sh{F}}:
  \sh{F}\longrightarrow
  \widehat{f}_*\bigl(\widehat{f}^{\,*}(\sh{F})\bigr)
\]
\sref[0]{0.4.4.3} are annihilated by
$\widehat{\sh{M}}^{\,n}$.
\end{lemma}

\begin{proof}
We may restrict to the case $X=\Spec(B)$, where $B$ is a
Noetherian ring; then $X'=\mathrm{V}(\mathfrak{K})$ for an ideal
$\mathfrak{K}$ of $B$.  We first reduce to the case in which $B$ is
adic Noetherian and $\mathfrak{K}$ is an ideal of definition.

Let $B_1$ be the separated completion of $B$ for the
$\mathfrak{K}$-preadic topology.
\oldpage[III]{155}
If $\mathfrak{K}_1=\mathfrak{K}B_1$, then $B_1$ is an adic
Noetherian ring and $\mathfrak{K}_1$ is an ideal of definition.
Put $X_1=\Spec(B_1)$, and let $h:X_1\to X$ be the morphism
corresponding to the canonical homomorphism $B\to B_1$.  If
$X'_1=h^{-1}(X')$, then $X'_1=\mathrm{V}(\mathfrak{K}_1)$.  Finally
put
\[
  Z_1=Z\times_X X_1,\qquad f_1:Z_1\longrightarrow X_1.
\]
The morphism $f_1$ is proper \sref[II]{II.5.4.2}.  Denote by
$\widehat{X}_1$ the completion of $X_1$ along $X'_1$, by
\[
  \widehat{Z}_1=
  Z_1\times_{X_1}\widehat{X}_1
\]
the completion of $Z_1$ along
$Z'_1=f_1^{-1}(X'_1)$, and by
$\widehat{f}_1$ the extension of $f_1$ to the completions.

The extension
$\widehat{h}:\widehat{X}_1\to\widehat{X}$ corresponds to the
identity map of $B_1$ and is therefore an isomorphism
\sref[I]{I.10.9.1}.  The corresponding morphism
$\widehat{Z}_1\to\widehat{Z}$ is consequently also an isomorphism,
and these isomorphisms identify $\widehat{f}_1$ with
$\widehat{f}$.  Moreover,
$\sh{M}_1=h^*(\sh{M})$ is a coherent ideal of $\sh{O}_{X_1}$ and
\[
  \Supp(\sh{O}_{X_1}/\sh{M}_1)
  =h^{-1}\bigl(\Supp(\sh{O}_X/\sh{M})\bigr)
\]
\sref[I]{I.9.1.13}.  Thus, if
$U_1=X_1-\Supp(\sh{O}_{X_1}/\sh{M}_1)$, then
$U_1=h^{-1}(U)$, and the restriction
$f_1^{-1}(U_1)\to U_1$ is an isomorphism
\sref[I]{I.3.2.7}.  Finally, $\widehat{\sh{M}}$ and
$\widehat{\sh{M}}_1$ are identified by $\widehat{h}$
\sref[I]{I.10.9.5}.  All the hypotheses of the lemma are therefore
satisfied by $X_1$, $X'_1$, $f_1$, and $\sh{M}_1$; henceforth we
may suppose that $B$ is adic Noetherian and that $\mathfrak{K}$ is
an ideal of definition.

We then have $\widehat{X}=\Spf(B)$ and
$\sh{F}=N^\Delta$, where $N$ is a finite-type $B$-module.  Thus
$\sh{F}=\widehat{\sh{G}}$, where $\sh{G}$ is the coherent
$\sh{O}_X$-module $\widetilde{N}$ \sref[I]{I.10.10.5}, and
\[
  \widehat{f}^{\,*}(\sh{F})
  =\widehat{f^*(\sh{G})}
\]
\sref[I]{I.10.9.5}.  By the first comparison theorem
\sref{III.4.1.5},
\[
  \widehat{f}_*
  \bigl(\widehat{f^*(\sh{G})}\bigr)
  \simeq
  \widehat{f_*\bigl(f^*(\sh{G})\bigr)},
\]
and, by \sref{III.5.1.3}, the canonical homomorphism
$\rho_{\sh{F}}$ is precisely $\widehat{\rho}_{\sh{G}}$.

The kernel $\sh{P}$ and cokernel $\sh{R}$ of
\[
  \rho_{\sh{G}}:
  \sh{G}\longrightarrow f_*\bigl(f^*(\sh{G})\bigr)
\]
are coherent $\sh{O}_X$-modules, and their restrictions to $U$ are
zero.  There is therefore an integer $n>0$ such that
\[
  \sh{M}^n\sh{P}=\sh{M}^n\sh{R}=0
\]
\sref[I]{I.9.3.4}.  It follows that
\[
  \widehat{\sh{M}}^{\,n}\widehat{\sh{P}}
  =
  \widehat{\sh{M}}^{\,n}\widehat{\sh{R}}
  =0
\]
\sref[I]{I.10.8.8} \sref[I]{I.10.8.10}, which proves the lemma.
\end{proof}

\begin{env}[5.3.5]
\label{III.5.3.5}
\emph{End of the proof of the existence theorem.}
Retain the hypotheses of \sref{III.5.1.4}.  We use Noetherian
induction \sref[0]{0.2.2.2}, assuming the theorem known for every
closed subprescheme $T$ of $X$ whose underlying space is distinct
from $X$; here $\widehat{T}$ denotes the completion of $T$ along
$T\cap X'$.  We may suppose that $X$ is nonempty.

Since $f$ is separated and of finite type, Chow's lemma
\sref[II]{II.5.6.1} gives a $Y$-scheme $Z$ and a $Y$-morphism
$g:Z\to X$ such that the structural morphism $h:Z\to Y$ is
quasi-projective, $g$ is projective and surjective, and there is a
nonempty open subset $U$ of $X$ for which
$g^{-1}(U)\to U$ is an isomorphism.  Let $\sh{M}$ be a coherent
ideal of $\sh{O}_X$ defining a closed subprescheme with underlying
space $X-U$ \sref[I]{I.5.2.2}, and let $\sh{F}$ be a coherent
$\sh{O}_{\widehat{X}}$-module whose support $E$ is proper over $Y$.
Denote by $\widehat{Z}$ the completion of $Z$ along
$h^{-1}(Y')$, and by
$\widehat{g}:\widehat{Z}\to\widehat{X}$ the extension of $g$ to
the completions.

The module $\widehat{g}^{\,*}(\sh{F})$ is a coherent
$\sh{O}_{\widehat{Z}}$-module whose support is contained in
$g^{-1}(E)$ and is therefore proper over $Y$, since $g$ is
projective and hence proper \sref[II]{II.5.4.6}.  Since $h$ is
quasi-projective, $\widehat{g}^{\,*}(\sh{F})$ is algebraizable by
\sref{III.5.2.6}.
\oldpage[III]{156}
It follows from \sref{III.5.3.3}, since $g$ is proper, that
\[
  \widehat{g}_*
  \bigl(\widehat{g}^{\,*}(\sh{F})\bigr)
\]
is an algebraizable $\sh{O}_{\widehat{X}}$-module.

Apply \sref{III.5.3.4} to $\sh{F}$ and $g$.  The kernel $\sh{P}$
and cokernel $\sh{R}$ of
\[
  \rho_{\sh{F}}:
  \sh{F}\longrightarrow
  \widehat{g}_*\bigl(\widehat{g}^{\,*}(\sh{F})\bigr)
\]
are annihilated by a power $\widehat{\sh{M}}^{\,n}$.  Let $T$ be
the closed subprescheme of $X$ defined by $\sh{M}^n$, whose
underlying space is $X-U$, and let $j:T\to X$ be the canonical
injection.  Its extension
$\widehat{j}:\widehat{T}\to\widehat{X}$ is the canonical injection
\sref[I]{I.10.14.7}, and we may write
\[
  \sh{P}=
  \widehat{j}_*\bigl(\widehat{j}^{\,*}(\sh{P})\bigr),
  \qquad
  \sh{R}=
  \widehat{j}_*\bigl(\widehat{j}^{\,*}(\sh{R})\bigr).
\]
Since $U$ is nonempty, the induction hypothesis shows that
$\widehat{j}^{\,*}(\sh{P})$ and
$\widehat{j}^{\,*}(\sh{R})$ are algebraizable
$\sh{O}_{\widehat{T}}$-modules.  Lemma \sref{III.5.3.3} then shows
that $\sh{P}$ and $\sh{R}$ are algebraizable, and
\sref{III.5.3.2} finally proves that $\sh{F}$ is algebraizable.
\end{env}


\subsection{Application: comparison between morphisms of ordinary preschemes and morphisms of formal preschemes. Algebraizable formal preschemes}
\label{subsection:III.5.4}

\begin{theorem}[5.4.1]
\label{III.5.4.1}
Let $A$ be an adic Noetherian ring, let $\mathfrak{J}$ be an ideal
of definition of $A$, and put
\[
  S=\Spec(A),\qquad S'=\mathrm{V}(\mathfrak{J}).
\]
Let $u:X\to S$ be a proper morphism and let $v:Y\to S$ be a
separated morphism of finite type.  Let
$\widehat{S}$, $\widehat{X}$, and $\widehat{Y}$ be the completions
of $S$, $X$, and $Y$ along $S'$, $u^{-1}(S')$, and $v^{-1}(S')$,
respectively.  For every $S$-morphism $f:X\to Y$, let
$\widehat{f}:\widehat{X}\to\widehat{Y}$ be its extension to the
completions.  Then the map $f\mapsto\widehat{f}$ is a bijection
\[
  \Hom_S(X,Y)
  \xrightarrow{\ \sim\ }
  \Hom_{\widehat{S}}(\widehat{X},\widehat{Y}).
\]
\end{theorem}

\begin{proof}
We first prove injectivity.  Suppose that $f,g:X\to Y$ are
$S$-morphisms such that $\widehat{f}=\widehat{g}$.  There is then
an open neighbourhood $V$ of $X'=u^{-1}(S')$ on which $f$ and $g$
coincide \sref[I]{I.10.9.4}.  Since $u$ is closed,
\sref{III.5.1.3.1} gives $V=X$, and hence $f=g$.

We now prove surjectivity.  Let
$h:\widehat{X}\to\widehat{Y}$ be an
$\widehat{S}$-morphism.  Put
\[
  Z=X\times_S Y,
\]
and denote the canonical projections by $p:Z\to X$ and
$q:Z\to Y$.  The prescheme $Z$ is of finite type over $S$
\sref[I]{I.6.3.4}, and is therefore Noetherian.  Let
$\widehat{Z}$ be its completion along
\[
  Z'=p^{-1}\bigl(u^{-1}(S')\bigr).
\]
The formal prescheme $\widehat{Z}$ is canonically identified with
$\widehat{X}\times_{\widehat{S}}\widehat{Y}$, and its projections
onto $\widehat{X}$ and $\widehat{Y}$ are identified with
$\widehat{p}$ and $\widehat{q}$ \sref[I]{I.10.9.7}.

Since $Y$ is separated over $S$, $\widehat{Y}$ is separated over
$\widehat{S}$ \sref[I]{I.10.15.7}.  The graph morphism
\[
  \Gamma_h=(1_{\widehat{X}},h):
  \widehat{X}\longrightarrow\widehat{Z}
\]
is consequently a closed immersion \sref[I]{I.10.15.4}.  Let
$\mathfrak{T}$ be the closed formal subprescheme of
$\widehat{Z}$ associated with this immersion, and let
$j:\mathfrak{T}\to\widehat{Z}$ be the canonical injection.  Thus
\[
  \Gamma_h=j\circ w,
\]
where $w:\widehat{X}\to\mathfrak{T}$ is an isomorphism
\sref[I]{I.10.14.3} whose inverse is
$\widehat{p}\circ j$.  Moreover, $\mathfrak{T}$ is proper over
$\widehat{S}$ because $\widehat{X}$ is.  It follows from
\sref{III.5.1.8} that there is a closed subprescheme $T$ of $Z$
such that
\[
  \mathfrak{T}
  =\widehat{T}
  =T_{/(T\cap Z')},
\]
and that $j=\widehat{i}$, where $i:T\to Z$ is the canonical
injection \sref[I]{I.10.14.7}.

The morphism $p\circ i:T\to X$ is an isomorphism: its extension to
the completions is the isomorphism
$\widehat{p}\circ\widehat{i}$, and it suffices to apply
\sref{III.4.6.8}, noting as above that $S$ is the only
neighbourhood of $S'$ in $S$.
\oldpage[III]{157}
Let $g:X\to T$ be the inverse of $p\circ i$, and put
\[
  f=q\circ i\circ g.
\]
Then $f:X\to Y$ has graph $\Gamma_f=i\circ g$.  Since
$\widehat{g}$ is the inverse of
$\widehat{p\circ i}$, we have
\[
  \widehat{\Gamma_f}
  =\widehat{i}\circ\widehat{g}
  =j\circ w
  =\Gamma_h.
\]
But $\widehat{\Gamma_f}=\Gamma_{\widehat{f}}$
\sref[I]{I.10.9.8}; hence $h=\widehat{f}$, which completes the
proof.
\end{proof}

In categorical language, the functor
$X\mapsto\widehat{X}$ is therefore fully faithful
\sref[0]{0.8.1.6} from the category of preschemes proper over
$\Spec(A)$ to the category of formal preschemes proper over
$\Spf(A)$, for every adic Noetherian ring $A$.  It consequently
gives an equivalence between the first category and a full
subcategory of the second.  The objects of this subcategory will
henceforth be called \emph{algebraizable formal preschemes}.  For
such a formal prescheme $\mathfrak{X}$, there is an ordinary
prescheme $X$, proper over $\Spec(A)$ and determined up to unique
isomorphism, such that $\mathfrak{X}\simeq\widehat{X}$.

\begin{env}[5.4.2]
\label{III.5.4.2}
\emph{Scholium.}
With the notation of \sref{III.5.4.1}, put
\[
  S_n=\Spec(A/\mathfrak{J}^{n+1}),\qquad
  X_n=X\times_S S_n,\qquad
  Y_n=Y\times_S S_n.
\]
It follows from \sref{III.5.4.1} and \sref[I]{I.10.12.3} that
giving an $S$-morphism $f:X\to Y$ is equivalent to giving an adic
inductive $(S_n)$-system \sref[I]{I.10.12.2} of
$S_n$-morphisms $f_n:X_n\to Y_n$.
\end{env}

\begin{remark}[5.4.3]
\label{III.5.4.3}
Contrary to what the existence theorem \sref{III.5.1.6} might
suggest, there are formal preschemes proper over $\Spf(A)$ which
are \emph{not algebraizable}, just as there are compact analytic
spaces which do not arise from complex algebraic varieties.  We
shall encounter such preschemes later in the ``theory of moduli,''
which, when the base field is $\mathbb{C}$, deals precisely with
infinitesimal variations of the complex structure of a complete
algebraic variety; it is well known that such variations can give
rise to analytic varieties which are not algebraic.
\end{remark}

\begin{proposition}[5.4.4]
\label{III.5.4.4}
Let $A$ be an adic Noetherian ring, put
$\mathfrak{S}=\Spf(A)$, and let
\[
  g:\mathfrak{X}\to\mathfrak{S},
  \qquad
  h:\mathfrak{Y}\to\mathfrak{S}
\]
be two proper morphisms of formal preschemes.  Let
$f:\mathfrak{X}\to\mathfrak{Y}$ be an
$\mathfrak{S}$-morphism.  If $f$ is finite and
$\mathfrak{Y}$ is algebraizable, then $\mathfrak{X}$ is
algebraizable.
\end{proposition}

\begin{proof}
The hypotheses on $g$ and $h$ already imply that $f$ is proper
\sref{III.3.4.1}; for $f$ to be finite, it is enough that the fibre
$f^{-1}(y)$ be finite for every $y\in\mathfrak{Y}$
\sref[0]{0.4.8.11}.  By hypothesis,
\[
  \sh{B}=f_*(\sh{O}_{\mathfrak{X}})
\]
is a coherent $\sh{O}_{\mathfrak{Y}}$-algebra
\sref[0]{0.4.8.6}.  Write
$\mathfrak{Y}=\widehat{Y}$ and $h=\widehat{w}$, where
$w:Y\to\Spec(A)$ is a proper morphism of ordinary preschemes.  The
existence theorem gives a coherent $\sh{O}_Y$-algebra $\sh{C}$
such that $\sh{B}=\widehat{\sh{C}}$.  Put
\[
  X=\Spec(\sh{C}),
\]
and let $u:X\to Y$ be the structural morphism.  The definition of
$\mathfrak{X}$ from $\sh{B}$ \sref[0]{0.4.8.7} then shows that
$\mathfrak{X}$ is canonically isomorphic to $\widehat{X}$ and
that $f$ is identified with $\widehat{u}$; it is enough to check
this when $Y$ is affine.
\end{proof}

Proposition \sref{III.5.1.8} is a special case of
\sref{III.5.4.4}.

\begin{theorem}[5.4.5]
\label{III.5.4.5}
Let $A$ be an adic Noetherian ring, let $\mathfrak{J}$ be an ideal
of definition of $A$, put
\[
  S=\Spec(A),\qquad
  \mathfrak{S}=\widehat{S}=\Spf(A),
\]
and let $f:\mathfrak{X}\to\mathfrak{S}$ be a proper morphism of
formal preschemes.  Put
\[
  S_k=\Spec(A/\mathfrak{J}^{k+1}),\qquad
  X_k=\mathfrak{X}\times_{\mathfrak{S}}S_k,
\]
and, for every $\sh{O}_{\mathfrak{X}}$-module $\sh{F}$, put
\[
  \sh{F}_k
  =\sh{F}\otimes_{\sh{O}_{\mathfrak{X}}}\sh{O}_{X_k}
  =\sh{F}/\mathfrak{J}^{k+1}\sh{F}.
\]
\oldpage[III]{158}
Let $\sh{L}$ be an invertible
$\sh{O}_{\mathfrak{X}}$-module, and suppose that
\[
  \sh{L}_0=\sh{L}/\mathfrak{J}\sh{L}
\]
is an ample $\sh{O}_{X_0}$-module.  Then $\mathfrak{X}$ is
algebraizable.  Moreover, if $X$ is a proper $S$-prescheme such
that $\mathfrak{X}\simeq\widehat{X}$, there is an ample
$\sh{O}_X$-module $\sh{M}$ such that
$\sh{L}\simeq\widehat{\sh{M}}$; in particular, $X$ is projective
over $S$.
\end{theorem}

\begin{proof}
Apply \sref{III.5.2.3} with
$\sh{F}=\sh{O}_{\mathfrak{X}}$.  There is an integer $n_0$ such
that, for $n\geq n_0$, the canonical homomorphism
\[
  \Gamma(\mathfrak{X},\sh{L}^{\otimes n})
  \longrightarrow
  \Gamma(X_0,\sh{L}_0^{\otimes n})
\]
is surjective.  Choose $n\geq n_0$ large enough that
$\sh{L}_0^{\otimes n}$ is very ample for $S_0$
\sref[II]{II.4.5.10}.  Since
$f_0:X_0\to S_0$ is proper,
$\Gamma(X_0,\sh{L}_0^{\otimes n})$ is a finite-type $A$-module
\sref{III.3.2.1}.  There is consequently a finite-type
$A$-submodule
\[
  E\subset\Gamma(\mathfrak{X},\sh{L}^{\otimes n})
\]
whose image in $\Gamma(X_0,\sh{L}_0^{\otimes n})$ is the whole
module.

For every $k\geq0$, consider the homomorphism
\[
  u_k:
  E/\mathfrak{J}^{k+1}E
  \longrightarrow
  \Gamma(X_k,\sh{L}_k^{\otimes n})
\]
induced by the canonical injection
$E\to\Gamma(\mathfrak{X},\sh{L}^{\otimes n})$.  The module
$(f_k)_*(\sh{L}_k^{\otimes n})$ is quasi-coherent, and
\[
  \Gamma\bigl(S_k,(f_k)_*(\sh{L}_k^{\otimes n})\bigr)
  =
  \Gamma(X_k,\sh{L}_k^{\otimes n}).
\]
Thus $u_k$ defines a homomorphism
\[
  \widetilde{u}_k:
  \widetilde{E/\mathfrak{J}^{k+1}E}
  \longrightarrow
  (f_k)_*(\sh{L}_k^{\otimes n}),
\]
and hence, by adjunction, a homomorphism
\[
  \widetilde{u}_k^{\,\sharp}:
  f_k^*\bigl(\widetilde{E/\mathfrak{J}^{k+1}E}\bigr)
  \longrightarrow
  \sh{L}_k^{\otimes n}.
\]
Put
\[
  \sh{G}_k=
  f_k^*\bigl(\widetilde{E/\mathfrak{J}^{k+1}E}\bigr).
\]
We have
$\sh{G}_0=\sh{G}_k/\mathfrak{J}\sh{G}_k$
\sref[I]{I.9.1.5}, and
\[
  \widetilde{u}_0^{\,\sharp}:
  \sh{G}_0/\mathfrak{J}\sh{G}_0
  \longrightarrow
  \sh{L}_0^{\otimes n}/
  \mathfrak{J}\sh{L}_0^{\otimes n}
\]
is obtained from $\widetilde{u}_k^{\,\sharp}$ by passage to the
quotients.

By the definition of $E$, $\widetilde{u}_0^{\,\sharp}$ is the
canonical homomorphism
\[
  \sigma:
  f_0^*\bigl((f_0)_*(\sh{L}_0^{\otimes n})\bigr)
  \longrightarrow
  \sh{L}_0^{\otimes n}.
\]
Since $\sh{L}_0^{\otimes n}$ is very ample,
$\widetilde{u}_0^{\,\sharp}$ is surjective
\sref[II]{II.4.4.3}.  Nakayama's lemma then shows that every
$\widetilde{u}_k^{\,\sharp}$ is surjective.

Each $\widetilde{u}_k^{\,\sharp}$ therefore defines
\sref[II]{II.4.2.2} an $S_k$-morphism
\[
  g_k:X_k\longrightarrow
  P_k=\mathbf{P}(E/\mathfrak{J}^{k+1}E).
\]
For $h\leq k$,
\[
  P_h=P_k\times_{S_k}S_h
\]
by \sref[II]{II.4.1.3}.  The relations among the
$\widetilde{u}_k^{\,\sharp}$ and \sref[II]{II.4.2.10} show that
$(g_k)$ is an adic inductive $(S_k)$-system
\sref[I]{I.10.12.2}.  The $g_k$ consequently define an
$\mathfrak{S}$-morphism of formal preschemes
\[
  g:\mathfrak{X}\longrightarrow\mathfrak{P},
\]
where $\mathfrak{P}$ is the inductive limit of the system $(P_k)$,
or equivalently the completion $\widehat{P}$ of
$P=\mathbf{P}(E)$.  The very ampleness of
$\sh{L}_0^{\otimes n}$ implies that $g_0$ is a closed immersion
\sref[II]{II.4.4.3}.  Hence $g$ is a closed immersion of formal
preschemes \sref[0]{0.4.8.10}, and $\mathfrak{X}$ is algebraizable
by \sref{III.5.1.8}.

The existence theorem \sref{III.5.1.6} now shows that $\sh{L}$ is
the completion $\widehat{\sh{M}}$ of an invertible
$\sh{O}_X$-module $\sh{M}$.  Moreover,
$\sh{L}^{\otimes n}$ is the completion of
$\sh{M}^{\otimes n}$ \sref[I]{I.10.8.10}, and the homomorphisms
$\widetilde{u}_k^{\,\sharp}$ define a unique homomorphism
\[
  v:f^*(\widetilde{E})
  \longrightarrow
  \sh{M}^{\otimes n}
\]
\sref{III.5.1.7}.  Since the
$\widetilde{u}_k^{\,\sharp}$ are surjective,
$\widehat{v}$ is surjective \sref[I]{I.10.11.5}, and therefore so
is $v$ \sref{III.5.1.3}.  Let $r:X\to P$ be the morphism defined
by $v$ \sref[II]{II.4.2.2}.  Its extension to the completions is
$g$.  Since $g$ is a closed immersion, so is $r$, by
\sref{III.5.1.8} and \sref{III.5.4.1}.  It follows that
$\sh{M}^{\otimes n}$ is very ample \sref[II]{II.4.4.6}, and hence
that $\sh{M}$ is ample \sref[II]{II.4.5.10}.
\end{proof}

\begin{remark}[5.4.6]
\label{III.5.4.6}
Let $A$ be an adic Noetherian ring, put
$\mathfrak{S}=\Spf(A)$, and let
$f:\mathfrak{X}\to\mathfrak{S}$ be a proper morphism of formal
preschemes.  Let $\sh{N}$ be the coherent ideal of
$\sh{O}_{\mathfrak{X}}$ such that, for every affine formal open
subset $U$ of $\mathfrak{X}$,
\[
  \Gamma(U,\sh{N})
\]
is the nilradical of
$\Gamma(U,\sh{O}_{\mathfrak{X}})$.  The existence of this ideal
follows readily from \sref[I]{I.10.10.2} and from the fact that
every ring homomorphism $B\to C$ sends the nilradical of $B$ into
that of $C$.  Let $\mathfrak{X}'$ be the closed formal
subprescheme of $\mathfrak{X}$ defined by $\sh{N}$
\sref[I]{I.10.14.2}.  It would be interesting to know whether
algebraizability of $\mathfrak{X}'$ implies algebraizability of
$\mathfrak{X}$ itself.

One would doubtless arrive at a solution of this problem if one
knew how to classify, for example by invariants
\oldpage[III]{159}
of a cohomological nature, the \emph{extensions} of a structural
sheaf $\sh{O}_{\mathfrak{X}}$---for an ordinary or a formal
prescheme---by a square-zero ideal.  Equivalently, these are the
$\sh{O}_{\mathfrak{X}}$-algebras $\sh{A}$ such that
$\sh{O}_{\mathfrak{X}}\simeq\sh{A}/\sh{J}$, where $\sh{J}$ is a
square-zero ideal of $\sh{A}$.
\end{remark}


\subsection{A decomposition of certain preschemes}
\label{subsection:III.5.5}

\begin{proposition}[5.5.1]
\label{III.5.5.1}
Let $A$ be an adic Noetherian ring, let $\mathfrak{J}$ be an ideal
of definition of $A$, and put $Y=\Spec(A)$.  Let $f:X\to Y$ be a
separated morphism of finite type, and put
\[
  Y_0=\Spec(A/\mathfrak{J}),\qquad
  X_0=X\times_Y Y_0=f^{-1}(Y_0).
\]
Let $Z_0$ be an open subset of $X_0$ which is proper over $Y_0$.
Then there is an open and closed subset $Z$ of $X$, proper over
$Y$, such that
\[
  Z\cap X_0=Z_0.
\]
\end{proposition}

\begin{proof}
By hypothesis, there is an open subset $T$ of $X$ such that
$T\cap X_0=Z_0$.  Let $\widehat{T}$ be the completion along $Z_0$
of the prescheme induced by $X$ on $T$.  The support of
$\sh{O}_{\widehat{T}}$ is $Z_0$, which is proper over $Y_0$;
hence $\widehat{T}$ is proper over
$\widehat{Y}=\Spf(A)$ \sref{III.3.4.1}.

It follows from \sref{III.5.1.8} that there is a closed
subprescheme $Z$ of $T$, proper over $Y$, such that, if
$i:Z\to T$ is the canonical injection, then
\[
  \widehat{i}:\widehat{Z}\longrightarrow\widehat{T}
\]
is an isomorphism; here $\widehat{Z}$ is the completion of $Z$
along $Z_0$.  By \sref{III.4.6.8}, there is an open neighbourhood
$V$ of $Z_0$ in $T$ such that
$i^{-1}(V)\to V$ is an isomorphism.  But $i^{-1}(V)$ is a
neighbourhood of $Z_0$ in $Z$, and is therefore equal to $Z$ by
\sref{III.5.1.3.1}.  Thus $Z$ is open in $T$, and hence in $X$.
It is also closed in $X$, since $Z$ is proper over $Y$ and $X$ is
separated over $Y$.  This proves the proposition.
\end{proof}

\begin{corollary}[5.5.2]
\label{III.5.5.2}
If $X_0$ is proper over $Y_0$, then $X$ is the union of two
disjoint open subsets $Z$ and $Z'$ such that $Z$ is proper over
$Y$ and contains $X_0$.  Moreover, every closed subset $P$ of $X$
which is proper over $Y$ is contained in $Z$.
\end{corollary}

\begin{proof}
Apply \sref{III.5.5.1} with $Z_0=X_0$, and put $Z'=X-Z$.  For the
last assertion, $P\cap Z'$ is closed in $P$ and is therefore proper
over $Y$.  If it were nonempty, $f(P\cap Z')$ would be a nonempty
closed subset of $Y$, and hence would meet $Y_0$
\sref{III.5.1.3.1}.  This contradicts
$Z'\cap X_0=\varnothing$, and proves that $P\subset Z$.
\end{proof}
